\labsection{6}{Build a blockchain from first principles}{sec:lab06-blockchain}

\begin{labproject}{Implement hashing, Merkle roots, chaining, and proof of work --- then break it}
\textbf{Coverage.} Chapter~\ref{ch:blockchain}.\\
\textbf{Goal.} Build the structure Chapter~\ref{ch:blockchain} describes, in about 120 lines of standard-library Python, and then demonstrate for yourself why altering history is computationally hard.

\begin{labmode}
This is the one lab in the course where you build the technology rather than analyse a market. Everything uses Python's built-in \texttt{hashlib} --- no packages to install, no network access, no cryptocurrency involved.

Work through the parts in order. Each one adds a single property, and the final property --- immutability --- only emerges once all the earlier ones are in place.
\end{labmode}

\medskip\noindent\textbf{Part A --- Hashing, and what makes it useful.}\par

\begin{lstlisting}[style=mlpython]
import hashlib

def sha256(data: str) -> str:
    return hashlib.sha256(data.encode()).hexdigest()

print(sha256("Alice pays Bob 10"))
print(sha256("Alice pays Bob 11"))   # one character different
\end{lstlisting}

Run it. The two outputs share no visible structure despite differing by one character. Four properties matter, and the rest of the lab depends on each:

\begin{mltable}
\begin{tabularx}{\linewidth}{@{}>{\bfseries\leavevmode\color{MLNavy}}p{0.22\linewidth}X@{}}
\toprule
\rowcolor{MLPaleBlue!92}
Property & Why the blockchain needs it \\
\midrule
Deterministic & The same input always hashes identically, so any node can verify any block. \\
Fixed length & A megabyte of transactions and one character both produce 64 hex characters. \\
Avalanche & One changed bit changes roughly half the output bits, so tampering is never subtle. \\
One-way & Given a hash you cannot compute the input, so mining must proceed by trial. \\
\bottomrule
\end{tabularx}
\end{mltable}

\medskip\noindent\textbf{Part B --- The Merkle root.}\par
A block commits to every transaction it contains through one 64-character value, built by hashing transactions in pairs until one remains:

\begin{lstlisting}[style=mlpython]
def merkle_root(transactions: list[str]) -> str:
    if not transactions:
        return sha256("")

    level = [sha256(tx) for tx in transactions]

    while len(level) > 1:
        # An odd node is paired with itself so every level halves cleanly.
        if len(level) % 2 == 1:
            level.append(level[-1])
        level = [sha256(level[i] + level[i + 1]) for i in range(0, len(level), 2)]

    return level[0]
\end{lstlisting}

\begin{keyidea}
The Merkle root means a block header of fixed size commits to an unlimited number of transactions. Change any transaction and the root changes, so the header no longer matches its own contents.

It also enables a proof that scales logarithmically: to convince someone a transaction is in a block of 1{,}024, you supply 10 hashes rather than all 1{,}024. That is what lets a phone verify a payment without storing the chain.
\end{keyidea}

\medskip\noindent\textbf{Part C --- Chaining blocks.}\par
Each block header contains the hash of the block before it. That single field is what makes the structure a chain rather than a list:

\begin{lstlisting}[style=mlpython]
from dataclasses import dataclass, field
import time

@dataclass
class Block:
    index: int
    transactions: list[str]
    previous_hash: str
    nonce: int = 0
    timestamp: float = field(default_factory=time.time)

    def header(self) -> str:
        return (f"{self.index}{self.timestamp}{self.previous_hash}"
                f"{merkle_root(self.transactions)}{self.nonce}")

    def hash(self) -> str:
        return sha256(self.header())
\end{lstlisting}

\medskip\noindent\textbf{Part D --- Proof of work.}\par
Mining searches for a nonce that makes the block's hash start with a required number of zeros. There is no shortcut: because the hash is one-way, the only method is to try values.

\begin{lstlisting}[style=mlpython]
def mine(block: Block, difficulty: int) -> Block:
    target = "0" * difficulty
    while not block.hash().startswith(target):
        block.nonce += 1
    return block
\end{lstlisting}

\begin{tryit}
Time mining at difficulty 1 through 6 and record the attempts:

\begin{lstlisting}[style=mlpython]
for d in range(1, 7):
    b = Block(index=1, transactions=["a pays b 1"], previous_hash="0" * 64)
    start = time.time()
    mine(b, d)
    print(f"difficulty {d}: {b.nonce:>9,} attempts, {time.time() - start:6.3f}s")
\end{lstlisting}

Each extra zero multiplies the expected work by 16, because each hex digit has 16 possible values.

A single run is very noisy: the nonce you find is a draw from a geometric distribution, so one sample can land at 5$\times$ or 30$\times$ the previous row through luck alone. Average eight trials per difficulty and the ratio settles close to 16 --- the supplied \texttt{blockchain.py} does this, and reports roughly 15.2$\times$ and 15.8$\times$ for the last two steps.

Now connect it to Chapter~\ref{ch:blockchain}: Bitcoin targets one block every ten minutes and adjusts difficulty every 2{,}016 blocks to hold that interval as total mining power changes. Explain what would happen to block times if difficulty were fixed and half the miners left.
\end{tryit}

\medskip\noindent\textbf{Part E --- Break it.}\par
This is the part that makes immutability concrete. Build a valid five-block chain, then alter a transaction in block 2:

\begin{lstlisting}[style=mlpython]
def validate(chain: list[Block], difficulty: int) -> list[str]:
    """Collect every problem, rather than returning at the first one."""
    problems = []
    for i in range(1, len(chain)):
        current, previous = chain[i], chain[i - 1]

        if not current.hash().startswith("0" * difficulty):
            problems.append(f"block {i}: proof of work invalid")
        if current.previous_hash != previous.hash():
            problems.append(f"block {i}: previous_hash does not match block {i-1}")

    return problems

chain[2].transactions[0] = "carol->mallory 1000"
for problem in validate(chain, difficulty=4):
    print(problem)
\end{lstlisting}

Collecting \emph{all} failures rather than returning at the first one is what makes the cascade visible. One edit produces two distinct errors:

\begin{lstlisting}[style=mlterminal]
block 2: proof of work invalid
block 3: previous_hash does not match block 2
\end{lstlisting}

The altered Merkle root changed block 2's own hash, so its proof of work no longer holds --- and that same changed hash is what block 3 stored, so the link breaks as well.

\begin{keyidea}
Editing one transaction changes that block's Merkle root, which changes its hash, which breaks the \texttt{previous\_hash} stored in the next block --- and every block after it.

To make the tampered chain valid again you must re-mine block 2 and every subsequent block, while the honest network extends the real chain ahead of you. That is what ``immutable'' means here. Not that the data cannot be altered, but that altering it costs more than the alteration is worth.
\end{keyidea}

\begin{pitfall}
A blockchain guarantees that recorded data has not changed. It guarantees nothing about whether the data was true when recorded.

Write a fraudulent transaction into a block, mine it correctly, and it is now permanently, verifiably fraudulent. This is the limit that matters for the supply-chain and land-registry applications in Chapter~\ref{ch:blockchain}: the ledger secures the record, not the fact.
\end{pitfall}

\begin{tryit}
Apply the chapter's use-case test to one non-financial application --- land registry, pharmaceutical provenance, or academic credentials.

Ask in order: are there multiple parties who do not trust each other, is a trusted intermediary genuinely absent or unacceptable, and does the record need to be shared rather than merely held? If any answer is no, a database with an audit log is simpler, faster, and cheaper. Say which your case is, and be willing to conclude that blockchain is the wrong tool.
\end{tryit}
\end{labproject}

\begin{verification}
\begin{itemize}
  \item \texttt{merkle\_root} handles an odd number of transactions and the empty case without crashing.
  \item Mining timings are recorded for at least difficulty 1--5, with the ratio between rows discussed.
  \item \texttt{is\_valid} detects both a broken link and invalid proof of work, and reports which block failed.
  \item The tamper demonstration shows the failure cascading beyond the edited block.
\end{itemize}
\end{verification}

\begin{labdeliverable}
Submit working \texttt{blockchain.py}, your difficulty-versus-time table with the observed scaling factor, the tamper-detection output, and one page applying the use-case test to a non-financial application --- including an honest verdict on whether a conventional database would serve it better.
\end{labdeliverable}
